% !PRINTED none
% [part-title / section-opening page: no running-head page number visible in any of the three bands; centred "Zweiter Abschnitt." / "Abel'sche Gruppen." title with rule, then \S.\,9 heading, then body text begins]
\begin{center}
\gesperrt{Second Chapter.}

\gesperrt{Abelian Groups.}

\rule{0.3\linewidth}{0.4pt}

\S.\,9.

\gesperrt{Representation of Abelian groups by a basis.}
\end{center}

We have already, in \S.\,1, called those groups in which the commutative law holds for composition commutative or \gesperrt{Abelian} groups. In these groups, which are of particular importance in the applications, much simpler laws prevail than in the general groups, as we shall now see. For the composition of the elements in these groups the same rules hold as in the multiplication of numbers, and we accordingly call the composita of elements of such a group, as in multiplication, products and powers.

Here too we restrict ourselves to the consideration of finite groups.

From the definition of Abelian groups it follows that the commutative law also holds for the composition of the parts of such a group (\S.\,4). Every divisor of an Abelian group is itself an Abelian group, and is at the same time a normal divisor, so that in these groups we may omit the addition ``Normal''.

A group which consists only of the repetitions of one element, such as
\[ 1, \; A, \; A^2 \; \ldots \; A^{a-1}, \]
we have already earlier (Bd. I, \S.\,148) called a cyclic group, and we retain this designation here. If $a$ is the smallest positive number for which $A^a = 1$, and thus $a$ is the degree of the group, then $a$ is also called the degree of the element $A$. If $m$ is any exponent for which $A^m = 1$, then $m$ is a
%CONT>>
